\documentclass[11pt]{article}
% Shared preamble for the social optimum derivations. Compile from this folder:
%   pdflatex pigou.tex   (the network diagrams are read from ../diagrams)
\usepackage[T1]{fontenc}
\usepackage{lmodern}
\usepackage{microtype}
\usepackage[margin=2.5cm]{geometry}
\usepackage{amsmath,amssymb}
\usepackage{booktabs}
\usepackage{float}
\usepackage{parskip}
\usepackage[font=small,labelfont=bf]{caption}
\usepackage[hyphens]{url}
\usepackage{standalone}
% Shared TikZ styles for the network diagrams.
\usepackage{tikz}
\usetikzlibrary{arrows.meta,positioning}
\definecolor{nodeblue}{HTML}{2A78D6}
\definecolor{endfill}{HTML}{E3EEFB}
\definecolor{shortcut}{HTML}{D64545}
\tikzset{
  v/.style={circle, draw=nodeblue, line width=1pt, minimum size=9mm, inner sep=0pt, font=\small},
  od/.style={v, fill=endfill},
  e/.style={-{Latex[length=2.6mm]}, line width=0.8pt},
  free/.style={e, shortcut},
  c/.style={font=\footnotesize, fill=white, inner sep=1.5pt},
}


\newcommand{\SCopt}{SC_{\mathrm{opt}}}
% Repository file or folder, e.g. \file{results/grid/summary.csv}; breaks at slashes.
\DeclareRobustCommand{\file}[1]{\path{#1}}


\begin{document}

\section{Social Optimum: $3 \times 3$ Grid Network}

\subsection{Network}

\begin{figure}[H]
  \centering
  \includestandalone[width=0.75\textwidth]{../diagrams/grid}
  \caption{Grid network. $x$ is the number of agents on an edge in a timestep; $2 \to 3$ is free.}
\end{figure}

All agents travel $1 \to 9$, moving right or down, on one of six paths:

\begin{center}
\begin{tabular}{@{}llr@{}}
\toprule
Path & Nodes & Base cost $\sum a_e$ \\
\midrule
RRDD & $1 \to 2 \to 3 \to 6 \to 9$ & 3.11 \\
RDRD & $1 \to 2 \to 5 \to 6 \to 9$ & 4.08 \\
RDDR & $1 \to 2 \to 5 \to 8 \to 9$ & 4.14 \\
DRRD & $1 \to 4 \to 5 \to 6 \to 9$ & 4.18 \\
DRDR & $1 \to 4 \to 5 \to 8 \to 9$ & 4.24 \\
DDRR & $1 \to 4 \to 7 \to 8 \to 9$ & 4.31 \\
\bottomrule
\end{tabular}
\end{center}

\subsection{Derivation}

The node in row $r$, column $c$ is reached at step $r + c$ on every path, so timing plays no role. With
$x_e$ agents on edge $e$,
\[
  SC = \sum_e x_e\,(a_e + 0.02\,x_e).
\]
At the optimum every used path has the same marginal cost $\sum_{e \in P}(a_e + 0.04\,x_e)$, and no unused
path is cheaper. There is no simple closed form, so we solve it numerically: add agents one at a time
along the path with the cheapest marginal cost, rerouting earlier agents when that helps (successive
shortest paths). This is exact for convex costs.

\textbf{Data.} Edge costs and $N$ from \file{scenarios/grid_N20}, \file{scenarios/grid_N50} and
\file{scenarios/grid}.

\subsection{Results}

\begin{table}[H]
\centering
\begin{tabular}{@{}lrrr@{}}
\toprule
Path & $N = 20$ & $N = 50$ & $N = 80$ \\
\midrule
RRDD & 15 & 28 & 39 \\
RDRD & 0  & 2  & 7  \\
RDDR & 0  & 0  & 0  \\
DRRD & 0  & 0  & 0  \\
DRDR & 3  & 10 & 17 \\
DDRR & 2  & 10 & 17 \\
\midrule
$\SCopt$   & 83.01  & 256.58 & 481.58  \\
$\SCopt/N$ & 4.1505 & 5.1316 & 6.01975 \\
\bottomrule
\end{tabular}
\caption{An optimal routing for each $N$. Edge loads are unique; the split into paths need not be.}
\end{table}

\begin{itemize}
  \item RRDD, the only path through the free edge, carries 75\% of agents at $N = 20$, 56\% at $N = 50$
        and 49\% at $N = 80$. As its other edges fill up, the optimum spreads agents across the grid.
  \item $N = 20$: the middle edges $2 \to 5$ and $5 \to 6$ are not worth using.
  \item $\SCopt/N$ rises from 4.15 to 6.02 as $N$ grows from 20 to 80.
\end{itemize}

\end{document}
